\uuid{zJYs}
\titre{Primitives usuelles}
\auteur{hanani ; mzari}
\organisation{exo7}
\licence{CC BY-NC-SA 4.0}
\datecreate{2020-01-20}
\niveau{L1}
\module{Analyse}
\chapitre{Calcul d'intégrales}
\sousChapitre{Primitives diverses}
\difficulte{1}
\theme{Primitives, primitives usuelles}
\exo7qcmid{lille-2.primitives.002}

\contenu{
\begin{question}
\qclassification{127.11}
\qtags{motcle=primitives usuelles}
Parmi les affirmations suivantes, cocher celles qui sont vraies :
\begin{answers}
\bad{$\displaystyle F(x)=x^2+\mathrm{e}^{x^2}$ est une primitive de $\displaystyle f(x)=2x+\mathrm{e}^{2x}$ sur $\Rr$.}
\good{$\displaystyle F(x)=x^2+\frac{\mathrm{e}^{2x}}{2}$ est une primitive de $\displaystyle f(x)=2x+\mathrm{e}^{2x}$ sur $\Rr$.}
\bad{$\displaystyle F(x)=x^2\mathrm{e}^{x}$ est une primitive de $\displaystyle f(x)=2x\mathrm{e}^{x}$ sur $\Rr$.}
\good{$\displaystyle F(x)=(2x-2)\mathrm{e}^{x}$ est une primitive de $\displaystyle f(x)=2x\mathrm{e}^{x}$ sur $\Rr$.}
\end{answers}
\begin{explanations}
On calcule $F'(x)$. Il en découle que $\displaystyle x^2+\frac{\mathrm{e}^{2x}}{2}$ est une primitive de $\displaystyle 2x+\mathrm{e}^{2x}$ et que $\displaystyle (2x-2)\mathrm{e}^{x}$ est une primitive de $\displaystyle 2x\mathrm{e}^{x}$.
\end{explanations}
\end{question}
}
